% !TeX root = ../drafts/06-bus-interactions.tex
\section{Bus interactions}\label{sec:omc}

The machine runs three interactions on the bus, told apart by their domain separators (\S\ref{sec:m3}): the VM state, the read-only memory, and the public bytecode. The last two are lookups, handled by offline memory checking~\cite{BEGKN94}.

\subsection{State}\label{sec:state}

The state interaction carries the 2 registers, program counter and frame pointer: $\tup{\dsep{ST},\pc,\fp}$. Each instruction row pulls its current state and pushes its successor $\tup{\dsep{ST},\nxt{\pc},\nxt{\fp}}$; the bus is preseeded by pushing the initial state $\tup{\dsep{ST},\pc_0,\fp_0}$ and pulling the final state $\tup{\dsep{ST},\pc_{\mathrm{final}},\fp_{\mathrm{final}}}$ (boundary conditions).
\begin{proposition}[Balance implies the existance of a valid execution]\label{prop:walk}
Let $\mathcal R$ be a finite multiset of rows, each pulling a state $u(r)$ and pushing a state $v(r)$, and suppose the state channel balances:
\[
  \mset{v(r):r\in\mathcal R}\uplus\mset{(\pc_0,\fp_0)}
  \;=\;\mset{u(r):r\in\mathcal R}\uplus\mset{(\pc_{\mathrm{final}},\fp_{\mathrm{final}})} .
\]
Then $\mathcal R$ splits into rows $r_1,\dots,r_L$ with
\[
  u(r_1)=(\pc_0,\fp_0),\qquad v(r_i)=u(r_{i+1}),\qquad v(r_L)=(\pc_{\mathrm{final}},\fp_{\mathrm{final}}),
\]
and rows forming closed walks.
\end{proposition}

\begin{proof}
Build the directed multigraph $G$ on the states: one edge $u(r)\to v(r)$ per row, plus one edge $e$ from $(\pc_{\mathrm{final}},\fp_{\mathrm{final}})$ to $(\pc_0,\fp_0)$.

Fix a state $x$. Its in-degree in $G$ counts the rows with $v(r)=x$, plus $e$ if $x=(\pc_0,\fp_0)$, where $e$ ends. Its out-degree counts the rows with $u(r)=x$, plus $e$ if $x=(\pc_{\mathrm{final}},\fp_{\mathrm{final}})$, where $e$ starts. Those two counts are the multiplicities of $x$ on the two sides of the hypothesis, so they are equal: $G$ is balanced at every state.

A finite balanced multigraph is a union of edge-disjoint closed walks: follow unused edges from any vertex, which balance always permits on entering one, until a vertex repeats; remove the closed walk found and repeat. Deleting $e$ from the closed walk that contains it leaves $r_1,\dots,r_L$.
\end{proof}

Tolerating closed walks is what lets a table be filled to a power of two (\S\ref{sec:e2e-pad}).

\subsection{The lookup protocol}\label{sec:lookup}

We use a similar idea to~\cite[Protocol.~4.27]{DP23}.

\paragraph{Setting.} A \emph{lookup array} $\lut$ holds $2^{\kappa}$ entries, entry $\lut[i]$ being a tuple of $\ncomp$ elements of $\K$. Its \emph{components} $\lut_0,\dots,\lut_{\ncomp-1}$ are the multilinears of height $2^{\kappa}$ holding one coordinate of the entry each: they are either committed or public. A \emph{read} is a row of a table (\S\ref{sec:m3}) naming an address $a\in\K$ and a value $v\in\K^{\ncomp}$, and is correct when $a=\gen^{i}$ and $v=\lut[i]$ for some $i<2^{\kappa}$.

\paragraph{Tuple and columns.} The bus interaction uses a tuple of the form $\tup{\sep,\addr,\cnt,\val}$: a domain separator naming the array, the address, the count, then the entry in coordinates $3,\dots,\ncomp+2$. Every higher coordinate is zero, so $\ncomp\le m-3$ for the bus width $m$ (\S\ref{sec:m3}). Take every part of a read to be a committed column of the reading table: its address, its count, and the $\ncomp$ components of the value read. Committed on top of $\lut$ is one finalize-count column $\cntfin$ of height $2^{\kappa}$.

\paragraph{Flush rules.} With $A[i]$ the number of reads of address $\gen^{i}$:
\begin{itemize}
\item \textbf{Seed:} for each $i<2^{\kappa}$, \push\ $\tup{\sep,\gen^{i},1,\lut[i]}$. A boundary block owned by no table, its address coordinate the index column of \S\ref{sec:idxcol}.
\item \textbf{Read} of $v$ at $a$ with count $\cnt$: \pull\ $\tup{\sep,a,\cnt,v}$ and \push\ $\tup{\sep,a,\gen\cdot\cnt,v}$, both flushes of the reading row.
\item \textbf{Finalize:} for each $i$, \pull\ $\tup{\sep,\gen^{i},\cntfin[i],\lut[i]}$. An honest prover commits $\cntfin[i]=\gen^{A[i]}$, the count the address has reached.
\end{itemize}

\paragraph{Completeness.} An honest prover gives the $t$-th read of $\gen^{i}$ the count $\gen^{t}$ and commits $\cntfin[i]=\gen^{A[i]}$: both sides then carry the counts $\gen^{0},\dots,\gen^{A[i]}$ at that address, so the bus balances.

\begin{lemma}[Invariant count multisets]\label{lem:invcnt}
Let $\rcnts$ be a finite multiset over $\K$ with $\gen\cdot\rcnts=\rcnts$. If $|\rcnts|<2^{64}-1$, every element of $\rcnts$ is $0$.
\end{lemma}

\begin{proof}
Let $\mu:\K\to\mathbb{N}$ be the multiplicity function of $\rcnts$ (Definition~\ref{def:multiset}). Multiplication by $\gen$ is a bijection of $\K$ and $\gen\cdot\rcnts$ has multiplicity function $c\mapsto\mu(\gen^{-1}c)$, so the hypothesis reads $\mu(\gen c)=\mu(c)$ for all $c$. Every nonzero $c$ is a power of $\gen$, so iterating this gives one value $\mu_1$ on all of $\K^{\times}$, which has $2^{64}-1$ elements. Then $|\rcnts|=\mu(0)+\mu_1(2^{64}-1)<2^{64}-1$ forces $\mu_1=0$.
\end{proof}

\begin{theorem}[Lookup correctness]\label{thm:omc}
Let the flush rules above be the only flushes carrying $\lut$'s separator, and write $(a_j,c_j,v_j)$ for the address, count and value that read $j$ evaluates to. If the bus balances, every $c_j\neq0$, and there are fewer than $2^{64}-1$ reads, then every read is correct: $(a_j,v_j)=(\gen^{i},\lut[i])$ for some $i<2^{\kappa}$.
\end{theorem}

\begin{proof}
Fix a read pair $(a,v)$ that is \emph{not} an entry of the array, so $(a,v)\neq(\gen^{i},\lut[i])$ for every $i<2^{\kappa}$. Balance equates the multiset of pushed tuples with the multiset of pulled ones; keep in each only the tuples carrying $\lut$'s separator, address $a$ and entry $v$, and the two are still equal. A seed and a finalization sit at the pairs $(\gen^{i},\lut[i])$, different from $(a,v)$, so neither survives. A read at $(a,v)$ does, on both sides at once: it pulls its count $c$ and pushes $\gen c$. So the survivors are exactly the reads at $(a,v)$.

Writing $\rcnts$ for the multiset of counts of the reads at $(a,v)$, balance therefore says $\gen\cdot\rcnts=\rcnts$. There are fewer than $2^{64}-1$ of them, so Lemma~\ref{lem:invcnt} makes every count in $\rcnts$ zero; every count is nonzero by hypothesis, so $\rcnts$ is empty. No read happens at such an invalid $(a,v)$.

\end{proof}

\begin{corollary}[Addresses come out in range]\label{cor:addrrange}
Under the same hypotheses every read address is an entry address $\gen^{i}$, $i<2^{\kappa}$: no read at any other address, $0$ included, occurs in an accepted proof.
\end{corollary}

\paragraph{Nothing checks the finalize counts.} $\cntfin$ appears nowhere in Theorem~\ref{thm:omc} or its proof. A wrong value (such as $0$ for instance) cannot make a bad read pass; it can only unbalance the bus at that entry, making the proof invalid.

\paragraph{Meeting the hypotheses.} Grand-product GKR checks balance (\S\ref{sec:gp}). It remains to check that all counts are nonzero, and the $\nread<2^{64}-1$ read bound, as follows.

\paragraph{The count product.} All counts are nonzero exactly when their product is. The prover therefore sends that product, one $\E$-element, and the verifier checks it is nonzero. It is proven like any other grand product, by GKR over the stacking of every count column (\S\ref{sec:leafstack}), batched with the two other GKRs of the bus (push / pull) (\S\ref{sec:gkr}).

\paragraph{Instance caps.} The verifier checks the announced sizes against public caps (\S\ref{sec:e2e-unrolled}): at most $2^{32}$ rows per table, at most $2^{32}$ memory cells (\S\ref{sec:memchan}) and at most $2^{32}$ bytecode instructions (\S\ref{sec:bytecode}). A row reads at most nine times, so the six tables flush $\nread\le9\cdot6\cdot2^{32}<2^{38}$ reads, far under the $2^{64}-1$ that Lemma~\ref{lem:invcnt} needs.

\paragraph{Relexing the hypotheses.} A bus entry may be any polynomial of degree at most $d$ in the row's columns (\S\ref{sec:m3}), so none of the three (address, count, value) need be committed. Example 1: we don't need to commit memory addresses of the form $\fp\cdot o$ (\S\ref{sec:memchan}) when $\fp$ and $o$ are already committed. Example 2: once the pull counter $\cnt$ is committed, no need to commit to the corrsponding push counter $\gen\cdot\cnt$.

\subsection{Memory}\label{sec:memchan}

The memory is an array of $2^{\kmem}$ cells, $16\le\kmem\le32$ (\S\ref{sec:vm}). Bus interaction uses domain separation $\sep=\dsep{MEM}$. A memory word is $192$-bit (i.e. one $\E$-element), so $\ncomp=3$, memory is committed as 3 $\K$-valued multilinears in $\kmem$ variables $M_0,M_1,M_2$.

\subsection{Bytecode}\label{sec:bytecode}

The bytecode is an array of $\nprog$ instructions, $\nprog\le2^{32}$ (\S\ref{sec:e2e-bc}). Bus interaction uses domain separation $\sep=\dsep{BC}$. An instruction is an opcode plus eight operand or immediate slots, so $\ncomp=9$: with the separator, the index and the counter ahead of them, this is the widest entry on the bus at $12$ coordinates, which is what forces the four index bits, hence the $m=16$ slots of \S\ref{sec:m3}. The program is public, i.e. not committed, the verifier can evaluate it directly.

\subsection{The index column}\label{sec:idxcol}

The seed and finalization blocks run over every address $\gen^{0},\dots,\gen^{2^{\kappa}-1}$, so the verifier needs the extension of that \emph{index column} $\mathrm{idx}$ at the point $\zeta$ (\S\ref{sec:leafstack}). It is not committed, and need not be: writing $w$ for the bits of $i$, $\gen^{\,i}=\prod_k(\gen^{\,2^{k}})^{w_k}$, whose $k$-th factor is $1+w_k(1+\gen^{\,2^{k}})$, which gives the efficient MLE formula:
\[
  \mle{\mathrm{idx}}(\zeta)\;=\;\prod_{k=0}^{\kappa-1}\bigl(1+\zeta_k\,(1+\gen^{\,2^{k}})\bigr),
\]
